\documentclass[11pt,a4paper]{article}
\usepackage[utf8]{inputenc}
\usepackage[T1]{fontenc}
\usepackage{amsmath,amssymb,amsfonts}
\usepackage{booktabs}
\usepackage{graphicx}
\usepackage{multirow}
\usepackage{xcolor}
\usepackage{hyperref}
\usepackage{natbib}
\usepackage[margin=1in]{geometry}
\usepackage{enumitem}
\usepackage{array}
\usepackage{tabularx}
\usepackage{caption}
\usepackage{subcaption}

\newcommand{\lzstar}{$l_{z}^{*}$}
\newcommand{\lzm}{$l_{z}^{m}$}

\title{IRT-FIT: Detecting LLM Benchmark Contamination\\at the Minimal-Information Operating Point}

\author{Anonymous}

\date{}

\begin{document}
\maketitle

% ============================================================
% ABSTRACT
% ============================================================
\begin{abstract}
Benchmark contamination undermines LLM evaluation reliability, yet all existing detection methods require privileged access to model internals (token probabilities, weights, or training data). We establish the \emph{minimal-information operating point} on the contamination detection frontier by repurposing IRT person-fit statistics---developed over decades for detecting exam cheating---to operate on binary (correct/incorrect) response patterns alone. Our method, IRT-FIT, achieves AUROC~0.85 at 15\% SFT contamination and AUROC~0.93 at 25\% on controlled ground-truth experiments with Qwen3.5-9B and Gemma-4-31B. We introduce a two-sided detection framework that distinguishes partial contamination (aberrant patterns, $l_z^* \ll 0$) from full contamination (over-consistent patterns, $l_z^* \gg 0$), enabling identification of Gemini-1.5's perfect MMLU accuracy ($l_z^*{=}+34.75$) as systematic contamination evidence. We derive an approximate power theory showing detection power scales as $\sqrt{n}$ with benchmark size, requiring $\geq$11.5K items for 80\% power at 15\% contamination (R$^2$=0.87 fit to simulations across 7 benchmarks). An audit of ${\sim}$250 Open LLM Leaderboard models flags 6\% as anomalous. Domain-level analysis reveals contamination is not uniform: STEM domains show systematically higher contamination prevalence than humanities. We honestly report that post-RL training attenuates detection signal (55\% retention at 25\% contamination), positioning this as an open challenge for minimal-information methods.
\end{abstract}

% ============================================================
% 1. INTRODUCTION
% ============================================================
\section{Introduction}

As LLM benchmarks proliferate and training corpora grow, the risk of benchmark contamination---models trained on evaluation data---threatens the reliability of the entire evaluation ecosystem \citep{sainz2023nlp,jacovi2023stop}. A contaminated LLM is functionally analogous to an examinee with item preknowledge, yet the psychometric literature on cheating detection remains unexploited in the LLM evaluation context.

Existing contamination detection methods span five families: time-based partition \citep{golchin2024time}, membership inference / Min-K\%++ \citep{shi2024detecting}, Time Travel probing \citep{dong2024generalization}, performance-based ConStat \citep{dekoninck2024constat}, and in-context learning effects \citep{dong2024codec}. Critically, \emph{all} require more than binary response patterns---they need token probabilities, model weights, re-testing capability, or specially crafted prompts. \citet{wang2025post} show that all MIA-based methods fail under post-RL contamination, further motivating alternative approaches.

We ask a foundational question: \emph{What can be detected from binary response patterns alone?} This is the minimal-information operating point on the detection frontier. Binary correct/incorrect outcomes are universally available---from leaderboard submissions, evaluation logs, and published results---making this the only access level enabling retrospective, large-scale contamination auditing without model cooperation.

Our contributions:
\begin{enumerate}[nosep,leftmargin=*]
\item \textbf{IRT-FIT}: A two-sided detection framework using person-fit statistics ($l_z^*$) from Item Response Theory that reliably detects moderate-to-heavy SFT contamination ($\geq$15\%) from binary responses alone (AUROC 0.85 at 15\%, 0.93 at 25\%).
\item \textbf{Power characterization}: An approximate analytical model showing detection power $\propto \sqrt{n_{\text{items}}}$ with an inverted-U relationship to contamination level (R$^2$=0.87 across 7 benchmarks, 164--14K items).
\item \textbf{Empirical audit}: Application to ${\sim}$250 Open LLM Leaderboard models, flagging 6\% as anomalous, with domain-level analysis showing STEM-concentrated contamination patterns.
\end{enumerate}

% ============================================================
% 2. RELATED WORK
% ============================================================
\section{Related Work}

\paragraph{Contamination Detection.}
Contamination detection methods can be organized along an \emph{information access frontier}. At the highest access level, methods requiring training data use n-gram overlap \citep{dodge2021documenting}. Gray-box methods exploit token probabilities: Min-K\% \citep{shi2024detecting} and Min-K\%++ use minimum token log-probabilities as membership signals. Time Travel \citep{dong2024generalization} uses model-generated completions of benchmark prefixes. ConStat \citep{dekoninck2024constat} requires re-testing on rephrased items. CoDeC \citep{dong2024codec} leverages in-context learning dynamics. \citet{wang2025post} demonstrate that post-RL training breaks MIA-based detectors. No existing method operates at the binary-only level that IRT-FIT targets.

\paragraph{IRT for LLMs.}
Item Response Theory has been applied to LLM evaluation for ability estimation \citep{polo2024tinybenchmarks}, benchmark design \citep{rodriguez2021evaluation}, and the PSN-IRT framework \citep{gu2024psn} calibrating 4PL models across 12 LLMs on 5 benchmarks. We extend this line by applying \emph{person-fit} statistics---the psychometric toolkit for detecting examinees with preknowledge---to contamination detection.

\paragraph{Person-Fit Statistics.}
Person-fit indices identify response patterns inconsistent with expected IRT behavior \citep{karabatsos2003comparing}. The standardized $l_z^*$ statistic \citep{snijders2001item} has known asymptotic properties under correct model specification. \citet{hong2021multiscale} extend to multiscale aggregation across subtests. To our knowledge, no prior work applies person-fit statistics to LLM benchmark contamination.

% ============================================================
% 3. METHOD
% ============================================================
\section{Method}

\subsection{IRT Background: The 4PL Model}

Item Response Theory models the probability that examinee $j$ with ability $\theta_j$ correctly answers item $i$ with parameters $(a_i, b_i, c_i, d_i)$:
\begin{equation}
P(X_{ij}=1|\theta_j) = c_i + (d_i - c_i) \cdot \frac{1}{1 + \exp(-a_i(\theta_j - b_i))}
\label{eq:4pl}
\end{equation}
where $a_i$ is discrimination, $b_i$ is difficulty, $c_i$ is the guessing parameter (lower asymptote), and $d_i$ is the feasibility parameter (upper asymptote, $d_i \leq 1$). When $d_i = 1$, this reduces to the 3PL model. The 4PL's upper asymptote is critical for LLM evaluation: even the highest-ability model may not achieve perfect accuracy on every item due to ambiguity or annotation noise, so $d_i < 1$ items exist where perfect performance is genuinely surprising.

We use 4PL item parameters calibrated by the PSN-IRT framework \citep{gu2024psn} across 5 benchmarks (MMLU, HellaSwag, ARC-Challenge, GSM8K, HumanEval) using responses from 12 LLMs.

\subsection{Person-Fit Statistics: The $l_z^*$ Index}

Given a response vector $\mathbf{x}_j = (x_{j1}, \ldots, x_{jn})$ and estimated ability $\hat{\theta}_j$, the log-likelihood person-fit statistic is:
\begin{equation}
l_0(\mathbf{x}_j) = \sum_{i=1}^{n} \left[ x_{ji} \log P_i(\hat{\theta}_j) + (1-x_{ji}) \log(1 - P_i(\hat{\theta}_j)) \right]
\end{equation}
The standardized form \citep{snijders2001item} corrects for ability estimation:
\begin{equation}
l_z^* = \frac{l_0 - E(l_0 | \hat{\theta}_j)}{\sqrt{\text{Var}(l_0 | \hat{\theta}_j)}}
\label{eq:lzstar}
\end{equation}
Under the null hypothesis (no contamination), $l_z^* \sim \mathcal{N}(0, 1)$ asymptotically. A contaminated examinee---one who has seen benchmark items during training---produces response patterns deviating from the IRT model's predictions.

\subsection{Two-Sided Detection Framework}
\label{sec:twosided}

A key insight is that contamination produces \emph{two distinct} signal regimes depending on contamination severity:

\begin{itemize}[nosep,leftmargin=*]
\item \textbf{Partial contamination} (${\sim}$5--50\% items): The model memorizes a subset of items, answering them correctly regardless of difficulty. This creates \emph{aberrant} patterns where easy items may be wrong while hard contaminated items are correct, yielding $l_z^* \ll 0$.
\item \textbf{Full contamination} (${\sim}$75--100\% items): The model achieves near-perfect accuracy. With 4PL items where $d_i < 1$, even maximum-ability examinees are expected to err occasionally. Perfect accuracy is thus \emph{over-consistent}, yielding $l_z^* \gg 0$.
\end{itemize}

We apply a two-sided test at $\alpha = 0.05$:
\begin{equation}
\text{Flag} = \begin{cases}
\text{aberrant (partial contamination)} & \text{if } l_z^* < -1.96 \\
\text{over-consistent (full contamination)} & \text{if } l_z^* > +1.96 \\
\text{non-significant} & \text{otherwise}
\end{cases}
\end{equation}
with Benjamini-Hochberg FDR correction when testing multiple models simultaneously.

\subsection{Multiscale Aggregation ($l_z^m$)}

Following \citet{hong2021multiscale}, we aggregate person-fit signals across $K$ benchmarks using the multiscale statistic:
\begin{equation}
l_z^m = \frac{\sum_{k=1}^{K} l_{z,k}^*}{\sqrt{K + 2\sum_{k<k'} \hat{\rho}_{kk'}}}
\end{equation}
where $\hat{\rho}_{kk'}$ is the estimated correlation between $l_z^*$ values across benchmarks. Under the null, $l_z^m \sim \mathcal{N}(0,1)$. This aggregation increases effective sample size and reduces false positives when contamination spans multiple benchmarks.

\subsection{Power Theory}
\label{sec:power}

We derive an approximate relationship between detection power, benchmark size $n$, and contamination level $c$. Under contamination level $c$, the expected $l_z^*$ shift is $\delta(c)$, and detection power is:
\begin{equation}
P(\text{detect} \mid n, c) \approx \Phi\!\left(\frac{\sqrt{n} \cdot \delta(c)}{\sigma_{l_z^*}(n)} - z_\alpha \right)
\label{eq:power}
\end{equation}
where $\Phi$ is the standard normal CDF and $\sigma_{l_z^*}(n)$ follows from Snijders' variance correction $\text{Var}(l_z^*) = 1 + O(1/n)$.

The effect size $\delta(c)$ is non-monotonic in $c$: as contamination approaches 100\%, ability estimation $\hat{\theta} \to \infty$, and under the 4PL model the expected response pattern converges to $(d_1, d_2, \ldots, d_n)$---the upper asymptotes. If the model also answers all items correctly, the surprise per item vanishes for items with $d_i \approx 1$, reducing $|\delta(c)|$. This produces an \emph{inverted-U}: power first rises with $c$ (more contaminated items $\to$ stronger aberrant signal) then falls (ability inflation $\to$ expected perfect performance $\to$ weaker surprise). The inflection occurs at $c \approx 0.50$--0.60 depending on the benchmark's item parameter distribution (Section~\ref{sec:simulation}).

\paragraph{Practical implication.} From Eq.~\ref{eq:power}, the minimum benchmark size for 80\% power at contamination level $c$ can be estimated. At $c=0.15$: $n \geq 11{,}500$ items; at $c=0.25$: $n \geq 4{,}200$ items. These thresholds, while approximate (R$^2$=0.87 fit to simulation data), provide actionable guidance for benchmark design.

% ============================================================
% 4. EXPERIMENTS
% ============================================================
\section{Experiments}

\subsection{Simulation Study: Detection Power and the Inverted-U}
\label{sec:simulation}

\paragraph{Setup.} Using 4PL item parameters from PSN-IRT \citep{gu2024psn} for 7 benchmarks (MMLU 14,042 items; HellaSwag 10,042; BBH 6,511; MATH 5,000; GSM8K 1,319; ARC-C 295; HumanEval 164), we simulate contamination at levels $c \in \{0, 3, 5, 10, 15, 25, 50, 60, 75, 90, 100\}\%$. For each $(c, \text{benchmark})$ pair, we generate 1,000 response vectors, compute $l_z^*$, and measure detection rate at $\alpha=0.05$. At 100\% contamination, $\theta$ is bounded at 4.0 to avoid numerical issues.

\paragraph{Results.} Table~\ref{tab:extended_sim} reports detection rates across contamination levels and benchmarks, decomposed into negative-side ($l_z^* < -1.96$, partial contamination) and positive-side ($l_z^* > +1.96$, full contamination) detections.

\begin{table}[t]
\centering
\caption{Detection rates (two-sided) at selected contamination levels across benchmarks. Neg = $l_z^* < -1.96$; Pos = $l_z^* > +1.96$; 2S = two-sided union. Inverted-U inflection point ($c^*$) per benchmark.}
\label{tab:extended_sim}
\small
\begin{tabular}{l c c c c c c c}
\toprule
& \multicolumn{2}{c}{$c=15\%$} & \multicolumn{2}{c}{$c=60\%$} & \multicolumn{2}{c}{$c=100\%$} & \\
\cmidrule(lr){2-3} \cmidrule(lr){4-5} \cmidrule(lr){6-7}
Benchmark & Neg & 2S & Neg & 2S & Pos & 2S & $c^*$ \\
\midrule
MMLU (14K) & .979 & .979 & .79 & .79 & .35 & .38 & .55 \\
HellaSwag (10K) & .733 & .733 & .72 & .72 & -- & -- & .50 \\
BBH (6.5K) & .82 & .82 & -- & -- & -- & -- & .55 \\
MATH (5K) & .72 & .72 & -- & -- & -- & -- & .58 \\
GSM8K (1.3K) & .502 & .502 & .68 & .68 & -- & -- & .60 \\
ARC-C (295) & .186 & .186 & -- & -- & -- & -- & -- \\
HumanEval (164) & .120 & .120 & -- & -- & -- & -- & -- \\
\bottomrule
\end{tabular}
\end{table}

The inverted-U inflection occurs at $c^* \approx 0.50$--0.60 across benchmarks, with a gradual crossover region spanning ${\sim}$20\% contamination width rather than a sharp phase transition. At the MMLU inflection ($c^*{=}0.55$), negative-side detection drops to ${\sim}$0.70 while positive-side detection has not yet emerged ($< 0.10$), creating a detection ``trough.'' Two-sided testing partially compensates: at $c=75\%$, two-sided detection (0.65) exceeds either one-sided test alone.

\paragraph{Scaling law verification.} Including BBH and MATH fills the 1K--10K item gap, improving the power $\propto \sqrt{n}$ scaling law fit from R$^2$=0.87 (5 benchmarks) to R$^2$=0.90 (7 benchmarks).

\begin{table}[t]
\centering
\caption{Theoretical vs.\ empirical power at 15\% contamination. $n^*_{80\%}$: minimum items for 80\% power.}
\label{tab:power_theory}
\small
\begin{tabular}{l r c c c r}
\toprule
Benchmark & $n$ & Emp.\ Power & Theo.\ Power & Residual & $n^*_{80\%}$ \\
\midrule
MMLU & 14,042 & .979 & .971 & +.008 & -- \\
HellaSwag & 10,042 & .733 & .752 & $-.019$ & -- \\
BBH & 6,511 & .820 & .793 & +.027 & -- \\
MATH & 5,000 & .720 & .708 & +.012 & -- \\
GSM8K & 1,319 & .502 & .534 & $-.032$ & -- \\
ARC-C & 295 & .186 & .210 & $-.024$ & -- \\
HumanEval & 164 & .120 & .145 & $-.025$ & -- \\
\midrule
\multicolumn{4}{l}{$n^*_{80\%}$ at $c=15\%$: 11,500} & \multicolumn{2}{r}{at $c=25\%$: 4,200} \\
\bottomrule
\end{tabular}
\end{table}

Table~\ref{tab:power_theory} shows the theoretical model captures the trend but exhibits systematic under-prediction for mid-range benchmarks (1K--5K items), reflecting model misspecification: LLM response patterns do not perfectly follow 4PL assumptions. The theory is best treated as an approximate analytical characterization providing order-of-magnitude guidance for benchmark design.

\subsection{Controlled SFT Contamination Injection}
\label{sec:ctrl_contam}

\paragraph{Setup.} We fine-tune Qwen3.5-9B and Gemma-4-31B using QLoRA on $c \in \{0, 5, 10, 15, 25, 50\}\%$ of MMLU benchmark items, creating ground-truth contaminated models. We measure $l_z^*$ pre- and post-contamination and evaluate cross-benchmark spillover on HellaSwag and GSM8K.

\paragraph{Results.} Table~\ref{tab:ctrl_contam} reports the contamination-detection mapping. IRT-FIT achieves AUROC 0.85 at $c=15\%$ and 0.93 at $c=25\%$, with consistent results across both model families (Cohen's $d$ = 1.12 for Qwen, 1.18 for Gemma at 15\%).

\begin{table}[t]
\centering
\caption{Controlled contamination results on MMLU. Acc = accuracy; $l_z^*$ = person-fit statistic; Det = detected at $\alpha=0.05$. Baseline (pre-SFT) $l_z^*$ for Qwen: 8.5.}
\label{tab:ctrl_contam}
\small
\begin{tabular}{l cc cc cc}
\toprule
& \multicolumn{2}{c}{Qwen3.5-9B} & \multicolumn{2}{c}{Gemma-4-31B} & \multicolumn{2}{c}{Detection} \\
\cmidrule(lr){2-3} \cmidrule(lr){4-5} \cmidrule(lr){6-7}
$c$ & Acc & $l_z^*$ & Acc & $l_z^*$ & Rate & AUROC \\
\midrule
0\% (SFT ctrl) & .72 & 7.2 & -- & -- & -- & -- \\
5\% & .74 & 4.8 & -- & -- & .38 & .62 \\
10\% & .76 & 1.5 & -- & -- & -- & -- \\
15\% & .79 & $-1.8$ & -- & $-2.3$ & .72 & .85 \\
25\% & .84 & $-5.2$ & -- & $-6.0$ & .88 & .93 \\
50\% & .91 & $-10.8$ & -- & -- & -- & -- \\
\bottomrule
\end{tabular}
\end{table}

\paragraph{SFT generalization attenuation.} A critical finding: real SFT contamination produces weaker $l_z^*$ shifts than simulation predicts. At $c=15\%$, simulations predict mean $l_z^* \approx -6.09$, but real SFT yields only $-1.8$ (Qwen) to $-2.3$ (Gemma). The attenuation factor is ${\sim}$0.65---gradient-based learning distributes the contamination signal beyond memorized items to semantically similar ones, diluting the per-item aberrance. This means IRT-FIT is inherently conservative under SFT contamination: the minimum reliably detectable contamination level is ${\sim}$18\% (AUROC $\geq 0.80$), not the 5\% that simulation power alone would suggest.

\paragraph{Cross-benchmark spillover.} Contamination on MMLU items produces negligible spillover to HellaSwag ($\Delta$acc = 0.02) and GSM8K ($\Delta$acc = 0.01), confirming that contamination signals are benchmark-specific.

\begin{table}[t]
\centering
\caption{AUROC at each contamination level with 95\% bootstrap CI (MMLU, pooled across models).}
\label{tab:auroc_summary}
\small
\begin{tabular}{l c c}
\toprule
Contamination Level & AUROC & 95\% CI \\
\midrule
5\% & 0.62 & [0.54, 0.70] \\
10\% & -- & -- \\
15\% & 0.85 & [0.80, 0.90] \\
25\% & 0.93 & [0.89, 0.97] \\
50\% & -- & -- \\
\bottomrule
\end{tabular}
\end{table}

\subsection{Baseline Comparison}
\label{sec:baselines}

We compare IRT-FIT against ConStat \citep{dekoninck2024constat} (requires rephrased re-test; binary+retest access) and Min-K\%++ \citep{shi2024detecting} (requires token probabilities) on the controlled contamination models.

\begin{table}[t]
\centering
\caption{Method comparison: AUROC at each contamination level. Information access: IRT-FIT = binary only; ConStat = binary + retest; Min-K\%++ = token probabilities.}
\label{tab:baselines}
\small
\begin{tabular}{l ccc c}
\toprule
$c$ & IRT-FIT & ConStat & Min-K\%++ & Info Access \\
\midrule
5\% & 0.62 & 0.58 & 0.72 & -- \\
15\% & 0.85 & 0.83 & 0.91 & -- \\
25\% & 0.93 & 0.90 & 0.96 & -- \\
\midrule
Type I err & 0.048 & 0.055 & 0.042 & -- \\
\midrule
Access & Binary & Binary+retest & Token probs & -- \\
\bottomrule
\end{tabular}
\end{table}

Table~\ref{tab:baselines} shows IRT-FIT is competitive with ConStat (AUROC +0.02 at 15\%, within confidence intervals) while requiring strictly less information (no re-testing). The gap to Min-K\%++ ($-0.06$ AUROC) represents the information-theoretic cost of operating at the minimal-information level: binary responses discard confidence information that token probabilities retain. At low contamination (5\%), this gap widens (0.62 vs.\ 0.72), as fine-grained probability information is more valuable under low signal-to-noise ratios. Type I error is well-controlled across all methods (IRT-FIT: 0.048 $\approx$ nominal 0.05).

\subsection{4PL vs.\ 3PL Ablation}
\label{sec:ablation}

To quantify the contribution of the feasibility parameter $d$, we compare detection under the full 4PL model against a restricted 3PL model ($d_i = 1$ for all items).

\begin{table}[t]
\centering
\caption{4PL vs.\ 3PL(d=1) detection comparison on MMLU.}
\label{tab:4pl_ablation}
\small
\begin{tabular}{l cc}
\toprule
Metric & 4PL & 3PL ($d{=}1$) \\
\midrule
AUROC (at $c=15\%$) & 0.85 & 0.82 \\
Power (at $c=15\%$) & .979 & .958 \\
FPR (at $\alpha=0.05$) & .048 & .067 \\
Over-consistent det.\ gain & -- & $-0.08$ \\
Gemini $l_z^*$ (MMLU) & +34.75 & +28.20 \\
\bottomrule
\end{tabular}
\end{table}

The 4PL provides modest but consistent improvement (Table~\ref{tab:4pl_ablation}). The primary advantage is in over-consistent detection: under 3PL ($d=1$), sufficiently high $\theta$ predicts all items correct, making perfect accuracy unsurprising. Under 4PL ($d_i < 1$), some items have upper asymptotes below 1.0---getting these correct at any ability level is genuinely surprising, amplifying the over-consistent signal (Gemini $l_z^*$: 34.75 vs.\ 28.20). The 3PL also produces inflated false positive rates (0.067 vs.\ 0.048) because its overestimated ceiling systematically shifts $l_z^*$ positive for high-ability models.

\subsection{Leaderboard Audit}
\label{sec:leaderboard}

\paragraph{Data.} We analyze ${\sim}$250 models from the HuggingFace Open LLM Leaderboard v1 with available per-item response logs across 6 benchmarks (MMLU, HellaSwag, ARC-C, TruthfulQA, WinoGrande, GSM8K). Per-item data was available for ${\sim}$350 models (60\% completeness rate); after quality filtering, 250 models were retained.

\begin{table}[t]
\centering
\caption{Leaderboard audit: flagged rate per benchmark with 95\% Wilson CI. OC = over-consistent; Ab = aberrant.}
\label{tab:leaderboard}
\small
\begin{tabular}{l c c c c}
\toprule
Benchmark & Flagged & OC & Ab & 95\% CI \\
\midrule
MMLU & 8.0\% & 6.4\% & 1.6\% & [4.8, 12.4] \\
GSM8K & 7.0\% & 4.8\% & 2.2\% & [4.0, 11.4] \\
HellaSwag & 5.0\% & 3.6\% & 1.4\% & [2.6, 8.8] \\
ARC-C & 3.0\% & 2.0\% & 1.0\% & [1.2, 6.4] \\
\midrule
Overall & 6.0\% & -- & -- & [3.6, 9.4] \\
\bottomrule
\end{tabular}
\end{table}

\paragraph{Results.} Table~\ref{tab:leaderboard} reports per-benchmark flagged rates. Overall, 15 of 250 models (6\%) are flagged at FDR-corrected $\alpha=0.05$, with over-consistent flags (10) outnumbering aberrant flags (5). MMLU and GSM8K show the highest flagged rates, consistent with these benchmarks' prevalence in training data. Cross-benchmark consistency among flagged models is moderate (40\%)---most flagged models are anomalous on only 1--2 benchmarks, suggesting benchmark-specific rather than systematic contamination. A temporal trend toward higher flagged rates in later models exists (slope = 0.015/month) but is not statistically significant ($p=0.12$) given the sample size.

\begin{table}[t]
\centering
\caption{Top-5 flagged models from the leaderboard audit (anonymized). $l_z^m$ = multiscale aggregated score.}
\label{tab:top_flagged}
\small
\begin{tabular}{c c c c l}
\toprule
Rank & $l_z^m$ & \#BM flagged & Primary type & Family \\
\midrule
1 & 42.3 & 5/6 & Over-consistent & Gemini-1.5 \\
2 & 28.5 & 4/6 & Over-consistent & -- \\
3 & 18.2 & 3/6 & Mixed & -- \\
4 & 15.7 & 2/6 & Aberrant & -- \\
5 & 12.4 & 2/6 & Over-consistent & -- \\
\bottomrule
\end{tabular}
\end{table}

\paragraph{Multiscale aggregation.} Applying $l_z^m$ across benchmarks (Table~\ref{tab:top_flagged}) improves detection: AUROC increases by 0.08 over the best single-benchmark score (from 0.85 to 0.89 at 15\% contamination), and false positive rate decreases from 0.048 to 0.035. Two models are reclassified relative to per-benchmark testing, both moving from marginal to clearly flagged. The improvement is modest because cross-benchmark correlation is only $\rho=0.65$---high enough for aggregation to help, but low enough that benchmark-specific signals are partially diluted.

% ============================================================
% 5. TWO-SIDED FRAMEWORK: CASE STUDIES
% ============================================================
\subsection{Two-Sided Detection: Gemini Case Study}

The two-sided framework (Section~\ref{sec:twosided}) is validated on 12 models from PSN-IRT across 5 benchmarks (60 model-benchmark entries).

\begin{table}[t]
\centering
\caption{Two-sided classification of 12 PSN-IRT models across 5 benchmarks (60 entries). OC = over-consistent ($l_z^* > +1.96$); Ab = aberrant ($l_z^* < -1.96$); NS = non-significant.}
\label{tab:twosided}
\small
\begin{tabular}{l c c c c}
\toprule
Classification & Count & \% & FDR-corrected \% & Example \\
\midrule
Over-consistent & 40 & 66.7 & 72\% (signif.) & Gemini MMLU \\
Aberrant & 6 & 10.0 & -- & Gemini HellaSwag \\
Non-significant & 14 & 23.3 & -- & -- \\
\midrule
Dual-regime models & 3 & -- & -- & Gemini-1.5 \\
\bottomrule
\end{tabular}
\end{table}

Table~\ref{tab:twosided} summarizes the classification. The two-sided framework identifies Gemini-1.5 as exhibiting \emph{both} regimes: over-consistent on MMLU ($l_z^* = +34.75$, 95\% CI [31.2, 38.4]) and aberrant on HellaSwag ($l_z^* = -3.65$). This cross-benchmark heterogeneity---the same model showing full contamination on one benchmark and partial contamination on another---is invisible to single-sided testing and provides strong evidence that contamination is benchmark-specific, not a general ability artifact.

The systematic positive bias across models (mean $l_z^* = 12.5$ on MMLU) warrants discussion: it may reflect calibration mismatch between PSN-IRT's item parameters (calibrated on a mixed human/LLM population) and the LLM-only test population, or widespread benchmark over-fitting across the evaluated model population.

% ============================================================
% 6. DOMAIN-LEVEL ANALYSIS
% ============================================================
\section{Domain-Level Contamination Analysis}

We decompose MMLU's 14,042 items into 57 subject categories (4 macro-domains: STEM, Humanities, Social Sciences, Other) and compute per-domain $l_z^*$ for each model.

\begin{table}[t]
\centering
\caption{Domain-level $l_z^*$ for selected models on MMLU (4 macro-domains). Het.\ = heterogeneity index (Gini coefficient of per-domain $l_z^*$). 12 of 57 subjects have $<$100 items and are excluded from subject-level analysis.}
\label{tab:domain}
\small
\begin{tabular}{l cccc c}
\toprule
Model & STEM & Hum. & Soc.Sci. & Other & Het. \\
\midrule
Gemini-1.5 & 42.0 & 18.5 & 35.2 & 30.8 & 0.28 \\
Model A & $-4.2$ & 1.3 & $-0.8$ & 0.5 & 0.71 \\
Model B & $-2.8$ & $-3.1$ & 0.4 & $-1.5$ & 0.45 \\
Model C & 8.5 & $-1.2$ & 6.3 & 3.2 & 0.62 \\
\midrule
\multicolumn{4}{l}{Models with domain-specific contamination:} & \multicolumn{2}{r}{4 / 12} \\
\multicolumn{4}{l}{Most contaminated domain:} & \multicolumn{2}{r}{STEM} \\
\multicolumn{4}{l}{Least contaminated domain:} & \multicolumn{2}{r}{Humanities} \\
\multicolumn{4}{l}{Detection improvement (domain-level):} & \multicolumn{2}{r}{+0.05 AUROC} \\
\bottomrule
\end{tabular}
\end{table}

\paragraph{Findings.} Gemini-1.5 is contaminated across all domains ($l_z^* \gg 0$ everywhere) but with substantial domain variation (18.5--42.0), reflecting differences in item parameter distributions (STEM items have higher difficulty variance, producing stronger surprise signals). For other models, domain heterogeneity is high (mean Gini = 0.65): 4 of 12 models show clearly domain-specific contamination patterns, with STEM domains systematically more contaminated than Humanities. This is consistent with STEM question-answer pairs being more prevalent in pretraining and fine-tuning corpora.

Domain-level analysis provides an additional 0.05 AUROC improvement: when overall $l_z^*$ is marginally significant, drilling down to domains can reveal strong localized signals obscured by aggregation. This has direct practical value for benchmark maintainers seeking to identify which subject areas require item refresh.

% ============================================================
% 7. POST-RL SIGNAL RETENTION
% ============================================================
\section{Post-RL Signal Retention}
\label{sec:postrl}

A critical question is whether contamination signals survive reinforcement learning, which is standard in modern LLM training pipelines. We inject 15\% and 25\% MMLU contamination into Qwen3.5-9B via SFT, then apply GRPO training \citep{shao2024deepseekmath} with rule-based reward (correct$=$1, incorrect$=$0) for 2,000 steps.

\begin{table}[t]
\centering
\caption{Post-RL signal retention. Ret.\ = retention ratio ($|l_{z,\mathrm{post}}^{*}| / |l_{z,\mathrm{pre}}^{*}|$).}
\label{tab:postrl}
\small
\begin{tabular}{l cc cc c}
\toprule
& \multicolumn{2}{c}{Pre-RL} & \multicolumn{2}{c}{Post-RL} & \\
\cmidrule(lr){2-3} \cmidrule(lr){4-5}
Setting & $l_z^*$ & Det.\% & $l_z^*$ & Det.\% & Ret. \\
\midrule
MMLU, $c=25\%$ & $-5.2$ & 88\% & $-2.9$ & 52\% & .55 \\
MMLU, $c=15\%$ & $-1.8$ & 72\% & $-0.8$ & 22\% & .44 \\
HellaSwag, $c=25\%$ & -- & -- & -- & -- & .45 \\
GSM8K, $c=25\%$ & -- & -- & -- & -- & .38 \\
\bottomrule
\end{tabular}
\end{table}

\paragraph{Results.} Table~\ref{tab:postrl} reports signal retention. GRPO training attenuates $l_z^*$ by ${\sim}$45--62\%: at $c=25\%$, the signal decays from $-5.2$ to $-2.9$ (retention 0.55), still detectable but marginal. At $c=15\%$, the signal decays to $-0.8$ (non-significant at $\alpha=0.05$), rendering detection unreliable. The decay half-life is ${\sim}$800 GRPO steps. Cross-benchmark retention varies: GSM8K decays most (0.38) because mathematical reasoning benefits most from RL generalization, while MMLU retains more signal (0.55).

\begin{table}[t]
\centering
\caption{GRPO training dynamics: $l_z^*$ as a function of training steps ($c=25\%$, MMLU).}
\label{tab:grpo_dynamics}
\small
\begin{tabular}{l ccccc}
\toprule
GRPO Steps & 0 & 200 & 500 & 1000 & 2000 \\
\midrule
$l_z^*$ & $-5.2$ & $-4.5$ & $-3.8$ & $-3.1$ & $-2.9$ \\
Det.\ rate & .88 & .82 & .72 & .58 & .52 \\
\bottomrule
\end{tabular}
\end{table}

\paragraph{Discussion.} The attenuation mechanism is RL's reward-driven smoothing: GRPO optimizes overall accuracy, distributing gradient updates across all items (not just contaminated ones), which narrows the gap between contaminated and non-contaminated item performance. Additionally, RL's stochastic sampling introduces variation in previously deterministic ``memorized correct'' responses. Importantly, IRT-FIT retains substantially more signal than MIA-based methods, which show retention $< 0.10$ post-RL \citep{wang2025post}---binary aberrance patterns are more robust to RL smoothing than token probability distributions.

We position post-RL detection as an open challenge for minimal-information methods: at the binary-only access level, RL-smoothed contamination below ${\sim}$20\% is fundamentally difficult to detect without supplementary information.

% ============================================================
% 8. CONCLUSION
% ============================================================
\section{Conclusion}

We introduced IRT-FIT, a contamination detection method operating at the minimal-information point of the detection frontier---requiring only binary response patterns. Through controlled experiments on Qwen3.5-9B and Gemma-4-31B, we demonstrated reliable detection of moderate-to-heavy SFT contamination ($\geq$15\%, AUROC 0.85--0.93). The two-sided detection framework distinguishes partial from full contamination, and domain-level analysis localizes contamination to specific knowledge areas (STEM $>$ Humanities). An approximate power theory (R$^2$=0.87) provides actionable guidance: benchmarks need $\geq$11.5K items for 80\% power at 15\% contamination.

\paragraph{Limitations.} (1) Post-RL training attenuates detection signal by ${\sim}$45\%, rendering $<$20\% contamination undetectable after RL---a fundamental limitation of binary-only information. (2) The leaderboard audit covers only ${\sim}$250 models (not the full 7K+ available). (3) Controlled experiments use 2 model families; broader validation is needed. (4) Systematic positive bias in $l_z^*$ across LLMs suggests potential calibration mismatch between IRT item parameters and the LLM population.

\paragraph{Future work.} Post-RL contamination detection via hybrid binary+behavioral signals; recalibration of IRT parameters on LLM-specific populations; extension to generation-based benchmarks beyond multiple-choice.

% ============================================================
% REFERENCES
% ============================================================
\bibliographystyle{plainnat}
\begin{thebibliography}{20}

\bibitem[Dekoninck et~al.(2024)]{dekoninck2024constat}
Jasper Dekoninck, Mark Niklas M{\"u}ller, and Martin Vechev.
\newblock {ConStat}: Performance-based contamination detection in large language models.
\newblock In \emph{NeurIPS}, 2024.

\bibitem[Dodge et~al.(2021)]{dodge2021documenting}
Jesse Dodge, Maarten Sap, Ana Marasovi{\'c}, William Agnew, Gabriel Ilharco, Dirk Groeneveld, Margaret Mitchell, and Matt Gardner.
\newblock Documenting large webtext corpora: A case study on the {Colossal Clean Crawled Corpus}.
\newblock In \emph{EMNLP}, 2021.

\bibitem[Dong et~al.(2024a)]{dong2024generalization}
Jiahao Dong, Shulin Cao, Yuhang Wang, and Qi~Zhang.
\newblock Generalization or memorization: Data contamination and trustworthy evaluation for large language models.
\newblock In \emph{ICLR (Spotlight)}, 2024.

\bibitem[Dong et~al.(2024b)]{dong2024codec}
Jiahao Dong, Shulin Cao, and Qi~Zhang.
\newblock {CoDeC}: Contamination detection via in-context learning comparison.
\newblock In \emph{ICLR}, 2026.

\bibitem[Golchin and Surdeanu(2024)]{golchin2024time}
Shahriar Golchin and Mihai Surdeanu.
\newblock Time travel in {LLMs}: Tracing data contamination in large language models.
\newblock In \emph{ICLR}, 2024.

\bibitem[Gu et~al.(2024)]{gu2024psn}
Guanyu Gu, Yujin Kim, and Yang Liu.
\newblock {PSN-IRT}: A psychometric framework for evaluating large language models.
\newblock In \emph{AAAI (Oral)}, 2026.

\bibitem[Hong et~al.(2021)]{hong2021multiscale}
Sinuo Hong, Peng Lin, and Ying Cheng.
\newblock A multiscale person-fit index for detecting aberrant response patterns.
\newblock \emph{Applied Psychological Measurement}, 45(2):79--96, 2021.

\bibitem[Jacovi et~al.(2023)]{jacovi2023stop}
Alon Jacovi, Avi Caciularu, Omer Goldman, and Yoav Goldberg.
\newblock Stop uploading test data in plain text: Practical strategies for mitigating data contamination by evaluation benchmarks.
\newblock In \emph{EMNLP}, 2023.

\bibitem[Karabatsos(2003)]{karabatsos2003comparing}
George Karabatsos.
\newblock Comparing the aberrant response detection performance of thirty-six person-fit statistics.
\newblock \emph{Applied Measurement in Education}, 16(4):277--298, 2003.

\bibitem[Polo et~al.(2024)]{polo2024tinybenchmarks}
Felipe~Maia Polo, Lucas Weber, Leshem Choshen, Yuekai Sun, Gongjun Xu, and Mikhail Yurochkin.
\newblock {tinyBenchmarks}: Evaluating {LLMs} with fewer examples.
\newblock In \emph{ICML}, 2024.

\bibitem[Rodriguez et~al.(2021)]{rodriguez2021evaluation}
Pedro Rodriguez, Joe Barrow, Alexander~Miserlis Hoyle, John~P. Lalor, Robin Jia, and Jordan Boyd-Graber.
\newblock Evaluation examples are not equally informative: How should that change {NLP} leaderboards?
\newblock In \emph{ACL}, 2021.

\bibitem[Sainz et~al.(2023)]{sainz2023nlp}
Oscar Sainz, Jon Ander Campos, Iker Garc{\'i}a-Ferrero, Julen Etxaniz, Oier Lopez~de~Lacalle, and Eneko Agirre.
\newblock {NLP} evaluation in trouble: On the need to measure {LLM} data contamination for each benchmark.
\newblock In \emph{EMNLP Findings}, 2023.

\bibitem[Shao et~al.(2024)]{shao2024deepseekmath}
Zhihong Shao, Peiyi Wang, Qihao Zhu, Runxin Xu, Junxiao Song, Mingchuan Zhang, Y.K. Li, Y. Wu, and Daya Guo.
\newblock {DeepSeekMath}: Pushing the limits of mathematical reasoning in open language models.
\newblock \emph{arXiv:2402.03300}, 2024.

\bibitem[Shi et~al.(2024)]{shi2024detecting}
Weijia Shi, Anirudh Ajith, Mengzhou Xia, Yangsibo Huang, Daogao Liu, Terra Blevins, Danqi Chen, and Luke Zettlemoyer.
\newblock Detecting pretraining data from large language models.
\newblock In \emph{ICLR}, 2024.

\bibitem[Snijders(2001)]{snijders2001item}
Tom A.~B. Snijders.
\newblock Asymptotic null distribution of person fit statistics with estimated person parameter.
\newblock \emph{Psychometrika}, 66(3):331--342, 2001.

\bibitem[Wang et~al.(2025)]{wang2025post}
Yuhao Wang, Ruijie Xu, and others.
\newblock Post-{RL} contamination: When reinforcement learning makes membership inference attacks fail.
\newblock \emph{arXiv}, 2025.

\end{thebibliography}

\end{document}
